% r03_h6w_c: proof of Theorem rot:thm:H6W, part 3: the constant-source shift on one lattice, the rate, the frozen family is Cauchy, and the assembly.

\subsubsection*{The constant-source shift on one lattice}

\begin{lemma}[the constant-source shift of the effective action]\label{rot:lem:H6SD}
Let $h,h'$ be in the family of \Cref{rot:lem:H6conv} with $h-h'=u\mathbf 1$ constant. Then on $\supp\chi_r$, for every real $t$,
\begin{equation}\label{rot:eq:H6SD}
\big|S^{(r)}_k(tG,h)-S^{(r)}_k(tG,h')\big|\le C(1+r)\,|u|\,p(\delta_r)^{n_1}\delta_r^{-3}|T|+2C\delta_r^\sigma|T|,
\end{equation}
with $C=C(t,G,h_0,|T|)$ independent of $k$ and $u$ and locally uniform in $t$, $n_1$ the exponent of \Cref{rot:prop:H6R}(ii) enlarged to absorb $(1+k-r)^{p_0}$, and $\sigma$ the remainder exponent.
\end{lemma}
\begin{proof}
\emph{Where $h$ enters.} By \cite[(3.14)--(3.17)]{K1} the Gaussian parts of $S^{(r)}_k$ are defined at the mass $m^2$ without sources, and the sources sit in $E'_r$ only, through the vertices (3.20), (3.21); the $h$-vertex is linear in $h$ with the explicit coupling $\delta_r^2\le1$. So $S^{(r)}_k(tG,h)-S^{(r)}_k(tG,h')=[P^{(r)}_k(h)-P^{(r)}_k(h')]+[R(h)-R(h')]$, $P$ the Taylor polynomial of order $\bar n=13$ and $R$ the remainder.
\emph{The remainder.} $|R(h)|,|R(h')|\le C\delta_r^\sigma|T|$ on $\supp\chi_r$, uniformly over the source family (\cite[p.~665]{K1}); this is the second term, and no derivative of $R$ in $u$ is taken.
\emph{The polynomial.} $P^{(r)}_k(h)$ is a finite sum of connected renormalised graphs with $m\ge0$ $h$-vertices, each multilinear in its source insertions; by multilinearity $P(h)-P(h')$ is the sum over the graphs with $m\ge1$ of $\sum_{i=1}^m[H$ with the $i$-th $h$-vertex carrying the constant $u$ and the others carrying $h$ or $h']$. The value of each such graph on $\supp\chi_r$ is $\sum_{\{y\},\{z\}}E^{(r)}_k(H;\{y\},\{z\})\times(\text{external fields at }\{y\})$. By \Cref{rot:prop:H6R}(ii), $|E^{(r)}_k(H)|\le C(1+r)(1+k-r)^{p_0}\delta_r^{\varrho_H}p(\delta_r)^{n_1}e^{-\delta\dist}$ with $\varrho_H\ge0$ and $C$ linear in $|u|$ for the distinguished vertex (a constant source has zero derivatives), the factor $(1+r)$ being present only for the degree-$0$ members of the order-$e^2$ $h$-linear field-free family and harmless elsewhere; the external scalar fields obey $|\phi_r(y)|\le p(\delta_r)(\lambda\delta_r)^{-1/4}$ \cite[(3.2)]{K1}, with $\delta_r^{-s/4}$ absorbed by $\varrho_H\ge s/4$ for $s$ external scalar legs, and the external photon legs are bounded likewise by (3.2), so the product of fields and kernel is $\le C|u|(1+r)(1+k-r)^{p_0}p(\delta_r)^{n_1+s}$ per position; this is the whole content of King's remark that the sources ``may be bounded when making a perturbative expansion'' and ``do not affect the leading order positivity of the action'' (p.~677). The graph is connected and its kernel decays exponentially in the tree distance of all its vertices, so the sum over all positions is $\le C|T^{(r)}|=C|T|\delta_r^{-3}$: one free position, the others tied. Collecting the finitely many graphs, with $(1+k-r)^{p_0}\le C(1+\log\delta_r^{-1})^{p_0}$ absorbed by enlarging $n_1$, gives the first term of \eqref{rot:eq:H6SD}.
\end{proof}

The bound \eqref{rot:eq:H6SD} keeps no smallness from the explicit $\delta_r^2$ of the $h$-vertex, from the running couplings or from the positive degrees: it is a per-site bound $O((1+r)p^{n_1})$ times the site count $|T|\delta_r^{-3}$. The true size of the worst term, the tree-level term $u\delta_r^2\sum\eta^3|\phi^{(r)}|^2$, is $|u|p^2\lambda^{-1/2}\delta_r^{-3/2}|T|$ on $\supp\chi_r$: a constant scalar source is a bare-mass shift, and its effect is a running-mass term of the same order as the mass term already present, which does not tend to zero. The smallness comes from the rate below, not from the effective action.

\begin{lemma}[the rate and the choice of $\vr$]\label{rot:lem:H6rate}
With $u_k=\frac12(\Delta_T(0)-\Delta_T(\eps_k))$ of \Cref{rot:lem:H6conv}, $|u_k|\le C\eps_k=C5^{-k}$ (\Cref{rot:prop:CT}(b)), and
\[
\eta'_k\defeq C(1+r)|u_k|p(\delta_r)^{n_1}\delta_r^{-3}+2C\delta_r^\sigma\le C(1+k)\,5^{-k+3\lfloor\vr k\rfloor}\big(b_0(1+k\ln5)\big)^{\mathrm pn_1}+2C\,5^{-\sigma\lfloor\vr k\rfloor}\longrightarrow0,
\]
since $\vr<\frac19<\frac13$ (\Cref{rot:lem:r}): the polynomial factor $(1+k)^{1+\mathrm pn_1}$ is beaten by the geometric $5^{-(1-3\vr)k}$. Hence, on $\supp\chi_r$, $|S^{(r)}_k(tG,h_k)-S^{(r)}_k(tG,\hinf)|\le\eta'_k|T|$, and, exponentiating under the common $\chi_r$ (same lattice, same depth, same fields),
\[
I_k(t;h_k)=e^{\theta_k(t)|T|}I_k(t;\hinf),\quad|\theta_k(t)|\le\eta'_k,\qquad\text{so}\qquad
\frac{\rho_k(t;h_k)}{\rho_k(t;\hinf)}\in[e^{-2\eta'_k|T|},e^{2\eta'_k|T|}],\qquad
|\rho_k(t;h_k)-\rho_k(t;\hinf)|\le16e^{C_Gt^2}\eta'_k|T|\to0
\]
for $k$ so large that $2\eta'_k|T|\le1$, locally uniformly in $t$.
\end{lemma}
\begin{proof}
The first display is \eqref{rot:eq:H6SD} with $u=u_k$, $\delta_r=5^{-\lfloor\vr k\rfloor}$, $1+r\le1+k$ and $p(\delta_r)\le b_0(1+k\ln5)^{\mathrm p}$; both terms tend to $0$ exactly when $0<\vr<\frac13$, and every condition of \Cref{rot:lem:r} on $\vr$ is an upper bound, so its window $\vr<\frac19$ needs no change. The second: $I_k(t;h_k)=\int\chi_re^{-S(tG,\hinf)}e^{-[S(tG,h_k)-S(tG,\hinf)]}$ lies in $[e^{-\eta'_k|T|},e^{\eta'_k|T|}]I_k(t;\hinf)$, and the same at $t=0$; divide, and use $\rho_k(t;\hinf)\le4e^{C_Gt^2}$ (\Cref{rot:lem:H6stepA}) with $e^{x}-1\le2x$ for $0\le x\le1$.
\end{proof}

\subsubsection*{The frozen family is Cauchy}

\begin{lemma}[one step at the frozen wall]\label{rot:lem:H6B}
Fix $t\in\R$ and write $r=r(k)$. For $k\ge k_0$ (so that $\delta_r<1$, $b_k,\tilde b_{k+1}\le\frac12$, $2C\delta_r^\sigma|T|\le\ln2$ and $2\eta_k|T|\le1$), with $\tilde\rho_{k+1}(t)\defeq\rho^{(r+1)}_{k+1}(t;\hinf)$ the small-field ratio of $X_{k+1}$ at depth $r+1$ (which is $r(k+1)$ or $r(k+1)+1$),
\[
\epsilon(k)\defeq|s_k(t)-s_{k+1}(t)|\le16e^{C_Gt^2}|T|\big(2C\delta_r^\sigma+\eta_k\big)+(1+4e^{C_Gt^2})\big(b_k+\tilde b_{k+1}\big),
\]
where $b_k\le2e^{(C+C_1)|T|}e^{-p(\delta_r)^2}5^{(q+q_1)k}$ and $\tilde b_{k+1}\le2e^{(C+C_1)|T|}e^{-p(\delta_r)^2}5^{(q+q_1)(k+1)}$. Consequently $\sum_k\epsilon(k)<\infty$, $(s_k(t))_k$ is Cauchy, and $W^{\rm tor}(t)\defeq\lim_ks_k(t)$ exists, locally uniformly in $t$.
\end{lemma}
\begin{proof}
\emph{(B1) The Cauchy estimate at one step.} By \Cref{rot:prop:H6I3}, on $\supp\chi_r$: $S^{(r)}_k(tG,\hinf)=S^{(r+1)}_{k+1}(tG,\hinf)+d_k+R_k(A_r,\phi_r;t)$ with $|R_k|\le\eta_k|T|$ and $d_k$ independent of the fields and of $t$. Exponentiate and integrate against the common $\chi_r$: $I^{(r)}_k(t;\hinf)=e^{-d_k}I^{(r+1)}_{k+1}(t;\hinf)e^{\theta|T|}$ with $|\theta|\le\eta_k$, and the same at $t=0$ with the same $d_k$. Dividing,
\[
\rho_k(t;\hinf)=\tilde\rho_{k+1}(t)\,e^{\theta'|T|},\quad|\theta'|\le2\eta_k,\qquad\text{so}\qquad|\rho_k(t;\hinf)-\tilde\rho_{k+1}(t)|\le\tilde\rho_{k+1}(t)\big(e^{2\eta_k|T|}-1\big).
\]
\emph{(B2) The large-field terms at the two depths.} \Cref{rot:lem:H6stepA} uses only \Cref{rot:prop:H6S1} at the given lattice and depth and \Cref{rot:prop:H6S2} for that lattice, which is depth-free; both hold for $\Xk$ at depth $r$ and for $X_{k+1}$ at depth $r+1$ ($\delta_{r+1}(X_{k+1})=\delta_r<1$), with the stated bounds on $b_k$ and $\tilde b_{k+1}$. Hence $\rho_k(t;\hinf),\tilde\rho_{k+1}(t)\le4e^{C_Gt^2}$ and, by \eqref{rot:eq:H6stepA} at both lattices,
\[
|s_k(t)-\rho_k(t;\hinf)|\le4e^{C_Gt^2}(e^{2C\delta_r^\sigma|T|}-1)+(1+4e^{C_Gt^2})b_k,\qquad
|s_{k+1}(t)-\tilde\rho_{k+1}(t)|\le4e^{C_Gt^2}(e^{2C\delta_r^\sigma|T|}-1)+(1+4e^{C_Gt^2})\tilde b_{k+1}.
\]
\emph{(B3) One step.} By (B1), (B2) and the triangle inequality,
\[
\epsilon(k)\le4e^{C_Gt^2}\big[2(e^{2C\delta_r^\sigma|T|}-1)+(e^{2\eta_k|T|}-1)\big]+(1+4e^{C_Gt^2})(b_k+\tilde b_{k+1}),
\]
and $e^x-1\le2x$ for $0\le x\le1$ gives the claim.
\emph{(B4) Summability.} With $\delta_r=5^{-\lfloor\vr k\rfloor}\le5\cdot5^{-\vr k}$ and $r=k-\lfloor\vr k\rfloor\ge(1-\vr)k$: $\delta_r^{\sigma}\le5^\sigma5^{-\sigma\vr k}$; $5^{-\gamma'r}\delta_r^{-\beta'}\le5^{\beta'}5^{-[\gamma'-(\gamma'+\beta')\vr]k}$, geometric once $\vr<\gamma'/(\gamma'+\beta')$, which holds for $\gamma'<\gamma_*$ close to $\gamma_*$ and $\beta'\in(\frac72,5)$, because the window $\vr<\gamma/(\gamma+\beta)$ of \Cref{rot:lem:r}, with $\gamma\le\gamma_*$ and $\beta\ge5$, lies strictly inside $\vr<\gamma_*/(\gamma_*+\beta')$; and $b_k+\tilde b_{k+1}\le4e^{(C+C_1)|T|}5^{q+q_1}\exp[-b_0^2(1+\lfloor\vr k\rfloor\ln5)^{2\mathrm p}+(q+q_1)k\ln5]$, a quadratic-in-$k$ negative exponent against a linear positive one ($\mathrm p\ge1$, \Cref{rot:lem:r}), hence $\le Ce^{-ck^2+C'k}$. So $\sum_{k\ge k_0}\epsilon(k)<\infty$, and for $k_0\le k<k'$, $|s_k(t)-s_{k'}(t)|\le\sum_{j=k}^{k'-1}\epsilon(j)\le\sum_{j\ge k}\epsilon(j)\to0$. All constants are uniform for $t$ in compacts. No uniformity in the number of steps is used or claimed: a sequence with summable consecutive differences is Cauchy.
\end{proof}

This lemma alone is King's torus theorem \cite[Thm.~2.1]{K1} with the frozen wall $\hinf$ as his scalar source, read on the normalised ratio: the differences from King's argument \cite[(3.10)--(3.13)]{K1} are that the drift-free stability bound $|\ln Z|\le C|T|$ \cite[(2.23)]{K1}, false at a wall, is replaced by the polynomial envelope of \Cref{rot:prop:H6S1,rot:prop:H6S2}, that the field-free $h$-sector constant, which does not converge, is removed by the ratio, and that the uniformity in the number of steps is replaced by summation over single steps.

\subsubsection*{Assembly}

\begin{proof}[Proof of \Cref{rot:thm:H6W}]
Fix $t\in\R$. By \Cref{rot:lem:H6stepA} at the two walls and \Cref{rot:lem:H6rate},
\[
|r_k(t)-s_k(t)|\le|r_k(t)-\rho_k(t;h_k)|+|\rho_k(t;h_k)-\rho_k(t;\hinf)|+|\rho_k(t;\hinf)-s_k(t)|\longrightarrow0 ,
\]
and by \Cref{rot:lem:H6B} $s_k(t)\to W^{\rm tor}(t)$. Hence $\lim_kr_k(t)=\lim_k\langle e^{tF(G)}\rangle_{\fc,\Xk}$ exists and equals $W^{\rm tor}(t)$, for every real $t$; by \Cref{rot:lem:H6conv} this is the limit of the torus-convention ratios at the moving wall $h_k=h_0-\Delta_T(\eps_k)/2$, which coincide with $r_k(t)$ term by term, and of those at the frozen wall $\hinf=h_0-\Delta_T(0)/2$. By \Cref{rot:lem:H6A} the convergence is locally uniform on $\C$, the limit $\Wv_\fc(zG)$ is entire, and all moments of $F(G)$ converge. The limits at $h_0$ itself in the two conventions are not claimed equal, and nothing below uses that.

\emph{What the proof does not use.} No value of the constants $d_k$; no Wick scheme; no convergence of $\Ztor_k(t;h)$ or $\Ztor_k(0;h)$ separately (both drift like $\eps_k^{\pm q}$); no comparison of two lattices, no rotated lattice, no skew torus, no hybrid flow and no centred blocks; no change of mass; the bridge \Cref{rot:prop:CT} only through the exact identity (c), the boundedness of $\Delta_T$ and the rate $|\Delta_T(s)-\Delta_T(0)|\le Cs$ of (b); and the sign of the wall nowhere, when the lower bound is taken from \Cref{rot:thm:FS}. The statement is at fixed $\fc$, $G$ and $t$, with $\eps_k\to0$ along $\eps_k=5^{-k}$; nothing depends on the later infimum over configurations or on the wall-height index of the shell construction.
\end{proof}
